\documentclass[10pt,a4paper]{article}

% ---------- Seiten- und Schrifteinstellungen ----------
\usepackage[a4paper,top=1.20cm,bottom=1.20cm,left=1.55cm,right=1.55cm]{geometry}
\usepackage[T1]{fontenc}
\usepackage[utf8]{inputenc}
\usepackage[ngerman]{babel}
\usepackage[scaled=0.96]{helvet}
\renewcommand{\familydefault}{\sfdefault}
\usepackage{microtype}
\usepackage{graphicx}
\usepackage{tabularx}
\usepackage{enumitem}
\usepackage{xcolor}
\usepackage{titlesec}
\usepackage[hidelinks]{hyperref}
\usepackage{ragged2e}
\usepackage{fancyhdr}
\usepackage{lastpage}

% ---------- PDF-Metadaten ----------
\hypersetup{
  pdftitle={REPLACENAME - Lebenslauf Windenergie},
  pdfauthor={REPLACENAME},
  pdfsubject={Windturbinenlasten, aeroelastische Simulation, Strukturdynamik und Python-Automatisierung}
}

% ---------- Stil ----------
\definecolor{heading}{HTML}{202020}
\definecolor{muted}{HTML}{555555}
\definecolor{rulecolor}{HTML}{777777}

\pagestyle{fancy}
\fancyhf{}
\renewcommand{\headrulewidth}{0pt}
\renewcommand{\footrulewidth}{0pt}
\fancyfoot[L]{\scriptsize\color{muted}REPLACENAME}
\fancyfoot[R]{\scriptsize\color{muted}\thepage/\pageref{LastPage}}

\setlength{\parindent}{0pt}
\setlength{\parskip}{0pt}
\setlist[itemize]{leftmargin=1.18em,itemsep=2.0pt,topsep=2.5pt,parsep=0pt,partopsep=0pt}
\urlstyle{same}
\RaggedRight

\titleformat{\section}
  {\large\bfseries\color{heading}}
  {}{0pt}{}
  [\vspace{-0.22em}\color{rulecolor}\titlerule]
\titlespacing*{\section}{0pt}{0.55em}{0.3em}

% Einträge über die volle Breite: Inhalt links, Datum bzw. Angabe rechts.
\newcommand{\cvjob}[4]{%
  \begin{tabularx}{\linewidth}{@{}Xr@{}}
    \textbf{#1} & \textbf{#2}\\[1pt]
    \textit{#3} & \textit{#4}
  \end{tabularx}
  \vspace{0.05em}
}

\newcommand{\cvproject}[2]{%
  \begin{tabularx}{\linewidth}{@{}Xr@{}}
    \textbf{#1} & \textbf{#2}
  \end{tabularx}
  \vspace{0.08em}
}

\newcommand{\cvedu}[4]{%
  \noindent
  \textbf{#1}\hfill\mbox{\textbf{#2}}\par
  \vspace{2pt}
  \noindent
  \textit{#3}\hfill\mbox{\textit{#4}}\par
  \vspace{0.3em}
}
\hyphenpenalty=7000
\exhyphenpenalty=7000
\emergencystretch=1em
\renewcommand{\arraystretch}{1.1}
\hyphenation{SOLIDWORKS OpenFAST}

\begin{document}
\fontsize{11.25}{12.75}\selectfont

% ---------- Kopfbereich ----------
\begin{minipage}[t]{0.79\linewidth}
  \vspace{0pt}
  {\fontsize{21.5}{23.5}\selectfont\bfseries REPLACENAME}\\[2pt]
  {\fontsize{11.55}{13.2}\selectfont\bfseries Maschinenbauingenieur | WEA-Lastanalyse \& Windparkbewertung}\\[4pt]
  REPLACELOCATION \enspace|\enspace
  \href{tel:REPLACEPHONE}{REPLACEPHONE} \enspace|\enspace
  \href{mailto:REPLACEEMAIL}{REPLACEEMAIL}\\[1.5pt]
  \href{REPLACELINKEDIN}{REPLACELINKEDIN} \enspace|\enspace
  \href{REPLACEPORTFOLIO}{REPLACEPORTFOLIO}\\[1.5pt]
\end{minipage}
\hfill
\begin{minipage}[t]{0.17\linewidth}
  \vspace{0pt}
  \raggedleft
  \IfFileExists{candidate_photo.jpeg}{\includegraphics[width=2.35cm,height=3.00cm,keepaspectratio]{candidate_photo.jpeg}}{\rule{0pt}{3.00cm}\hspace{2.35cm}}
\end{minipage}

\section{Kurzprofil}
Maschinenbauingenieur mit nahezu sechs Jahren Industrieerfahrung in Produktentwicklung, Strukturdynamik und simulationsgestützter Konstruktion, mit aktueller Spezialisierung auf Windenergie. Erfahrung mit aeroelastischen OpenFAST-Simulationen, Auslegungslastfällen nach IEC 61400-1, Python-Automatisierung, Nachverarbeitung von Lastdaten sowie Windparkplanung in windPRO und QGIS. Verbindet Windenergie-Simulation mit praktischer Maschinen- und Systementwicklung bis zu Fertigung und Installation.

\section{Technische Kernkompetenzen}
\begin{tabularx}{\linewidth}
{@{}>{\raggedright\arraybackslash}p{4.5cm}X@{}}
\textbf{WEA-Simulation} & OpenFAST, aeroelastische Modellierung, Auslegungslastfälle nach IEC 61400-1, Analyse von Anlagensignalen und Strukturlasten\\[2pt]
\textbf{Engineering-Automatisierung} & Python-Skripte, Generierung von Simulationsfällen, Batch-Ausführung, Monitoring, Datenorganisation und Nachverarbeitung\\[2pt]
\textbf{Windparkplanung} & windPRO, QGIS, Anlagenlayout, Energieertragsbewertung, Schallimmissionen, Schattenwurf und Abschaltungen zum Fledermausschutz\\[2pt]
\textbf{Strukturdynamik} & ANSYS Workbench, Modalanalyse, harmonische Analyse, Schwingungs- und Resonanzbewertung\\[2pt]
\textbf{Maschinenbau} & CATIA V5, SOLIDWORKS, Siemens NX, fertigungsgerechte Konstruktion, Zeichnungen und Stücklisten\\[2pt]
\textbf{Programmierung/Data} & Python, MATLAB/Simulink, Power BI, Excel
\end{tabularx}

\section{Ausgewählte Erfahrung in der Windenergie}
\cvproject{Masterarbeit - Rotorbasierte Windzuflussschätzung}{Laufend}
\begin{itemize}
  \item Entwicklung eines Rotor-als-Sensor-Ansatzes zur Schätzung von Windgeschwindigkeit, Windrichtung, vertikaler Windscherung und Winddrehung aus messbaren, mit OpenFAST simulierten Anlagensignalen.

  \item Automatisierung der Simulations- und Signalverarbeitungsabläufe in Python sowie Vergleich interpretierbarer linearer Schätzer mit Modellen auf Basis neuronaler Netze.
\end{itemize}
\vspace{1pt}
\cvproject{20-MW-Offshore-Windenergieanlage - Aeroelastische Lastanalyse}{Hochschulprojekt}
\begin{itemize}
  \item Durchführung von OpenFAST-Simulationen für eine 20-MW-Offshore-Windenergieanlage mit Monopile-Gründung für Auslegungslastfälle nach IEC 61400-1 mit Python- und MATLAB-gestützten Abläufen.
  \item Nachverarbeitung von Anlagensignalen sowie Bewertung von Strukturlasten und dynamischem Verhalten unter unterschiedlichen Umwelt- und Betriebsbedingungen.
  \item Erstellung technischer Diagramme und vergleichender Ergebnisdarstellungen zur Interpretation lastbestimmender Bedingungen und des Anlagenverhaltens.
\end{itemize}
\vspace{1pt}
\cvproject
  {\href{REPLACEPROJECTLINK}
  {OpenFAST-Workflow-Automatisierung \& Output-Monitoring}}
  {Hochschulprojekt}

\begin{itemize}
  \item Entwicklung von Python-Werkzeugen für OpenFAST-Fallgenerierung, Batch-Ausführung und automatisierte Ablage der Ausgabedaten.

  \item Entwicklung eines Python-Tools zur Echtzeit- und Postprocessing-Visualisierung von OpenFAST-Ausgaben für die schnelle Prüfung von Betriebsparametern, Strukturlasten und dynamischer Anlagenreaktion.
\end{itemize}
\vspace{1pt}
\cvproject{Windparkplanung und Umweltbewertung}{Hochschulprojekt}
\begin{itemize}
  \item Durchführung einer standortbezogenen Windparkplanung in windPRO und QGIS, einschließlich Anlagenlayout, Energieertragsschätzung, Schallimmissionsbewertung und Schattenwurfanalyse.

  \item Bewertung fledermausbedingter Abschaltungen und ihres Einflusses auf den jährlichen Energieertrag.

  \item Erstellung von Video-Tutorials zu ausgewählten windPRO-Workflows und zur Ergebnisinterpretation.
\end{itemize}

\newpage

\section{Berufserfahrung}
\cvjob{Entwicklungsingenieur - Solar-Trackingsysteme \& automatisierte Mechanismen}{01/2022--02/2023}{Paydar | Solarenergiesysteme}{Teheran}
\begin{itemize}
  \item Konstruktion von Photovoltaik-Nachführsystemen, Befestigungskomponenten und mechanischen Baugruppen in SOLIDWORKS, unterstützt durch strukturelle Analysen in ANSYS.
  \item Entwicklung eines elektromotorisch angetriebenen Reinigungssystems für Solarmodule mit Schienen- und Radführung sowie Begleitung durch Detailkonstruktion, Fertigung, Montage und erfolgreiche Installation vor Ort.
  \item Überarbeitung mechanischer Komponenten zur Verbesserung der Herstellbarkeit; Beitrag zu einer Reduzierung von Produktionsausschuss und Rohmaterialabfall um etwa 10\%.
  \item Erstellung von Fertigungszeichnungen, Stücklisten und technischer Dokumentation in SAP-basierten Systemen.
\end{itemize}

\cvjob{Maschinenbauingenieur - Maschinenentwicklung \& Strukturdynamik}{03/2017--12/2021}{TehranRigzar | Bergbau- \& Aufbereitungstechnik}{Teheran}
\begin{itemize}
  \item Konstruktion mechanischer Anlagen und Baugruppen für die Gesteinsaufbereitung, darunter Backenbrecher, Förderbänder, Hochfrequenz-Schwingsiebe und Waschanlagen für Zuschlagstoffe.

  \item Erstellung von Maschinenbaugruppen, Strukturbauteilen, mechanischen Schnittstellen und Fertigungszeichnungen in CATIA V5.

  \item Mitentwicklung eines Hochfrequenz-Schwingsiebs auf Grundlage von Laboruntersuchungen zur Siebung und definierten Betriebsanforderungen.

  \item Durchführung von Modal- und harmonischen Analysen in ANSYS zur Bewertung des dynamischen Verhaltens des Siebsystems.

\end{itemize}

\section{Studium}

\cvedu
  {M.Sc. Wind Energy Engineering -- Notenschnitt: 1,9}
  {03/2023--heute}
  {Hochschule Flensburg, Flensburg, Deutschland}
  {Abschluss: Februar 2027}

\cvedu
  {M.Sc. Mechanical Engineering}
  {09/2014--01/2017}
  {Hakim Sabzevari University, Sabzevar}
  {Regelungstechnik und Robotik}

\vspace{0.6mm}

\textit{Abschlussarbeit:} Entwicklung einer MATLAB/Simulink-Bewegungssteuerung zur präzisen Positionierung eines autonomen Schweißroboters für Unterwassereinsätze.

\vspace{1.2mm}

\cvedu
  {B.Sc. Mechanical Engineering}
  {09/2008--06/2013}
  {Islamic Azad University, Teheran}
  {Maschinenkonstruktion und Schwingungsanalyse}
\section{Technische Tools}
\begin{tabularx}{\linewidth}{@{}>{\raggedright\arraybackslash}p{4.75cm}X@{}}
\textbf{Windenergie} & OpenFAST, windPRO, QGIS\\
\textbf{Programmierung \& Daten} & Python, MATLAB/Simulink, Power BI, Excel\\
\textbf{CAE} & ANSYS Workbench, Abaqus\\
\textbf{CAD} & CATIA V5, SOLIDWORKS, Siemens NX\\
\textbf{Engineering-Systeme} & SAP, MS Project, technische Dokumentation
\end{tabularx}

\section{Ausgewählte Weiterbildungen \& Zertifikate}
\begin{tabularx}{\linewidth}{@{}>{\raggedright\arraybackslash}p{4.75cm}X@{}}
\textbf{Python-Programmierung} & HarvardX\\
\textbf{ANSYS} & IEEE-Schulung\\
\textbf{SOLIDWORKS} & Technical and Vocational Training Organization\\
\textbf{CATIA V5} & Amirkabir University of Technology
\end{tabularx}

\section{Sprachen}
\begin{tabularx}{\linewidth}{@{}>{\raggedright\arraybackslash}p{4.75cm}X@{}}
\textbf{Englisch} & C1; IELTS-Gesamtergebnis 7,0\\
\textbf{Deutsch} & B1; Vorbereitung auf eine Zertifikatsprüfung
\end{tabularx}

\end{document}
